\documentclass[10pt,aspectratio=169]{beamer}

\usetheme{Madrid}
\usecolortheme{seahorse}
\setbeamertemplate{navigation symbols}{}
\setbeamertemplate{footline}[frame number]

%\usepackage[spanish,es-tabla]{babel}
%\usepackage[utf8]{inputenc}
%\usepackage[T1]{fontenc}
%\usepackage{lmodern}
%\usepackage{amsmath,amssymb}
%\usepackage{graphicx}
%\usepackage{booktabs}
\usepackage{tabularx}
%\usepackage{array}
%\usepackage{xcolor}
%\usepackage{listings}
%\usepackage{tikz}
%\usetikzlibrary{arrows.meta,positioning,shapes.geometric}

\usepackage[utf8]{inputenc}
\usepackage[T1]{fontenc}
\usepackage[spanish]{babel}
\usepackage{amsmath, amssymb, booktabs, array, multirow}
\usepackage{graphicx}
\usepackage{xcolor}
\usepackage{listings}
\usepackage{tikz}
\usetikzlibrary{positioning,arrows.meta}
\usepackage{pgfplots}
\pgfplotsset{compat=1.18}

\definecolor{codebg}{RGB}{248,248,248}
\definecolor{codegray}{RGB}{80,80,80}
\definecolor{darkblue}{RGB}{20,55,100}
\definecolor{darkgreen}{RGB}{0,90,60}

\definecolor{uohblue}{HTML}{1F4E79}
\definecolor{softgray}{HTML}{F2F4F7}
\definecolor{darkgray}{HTML}{344054}
\definecolor{codegreen}{HTML}{2E7D32}
\definecolor{codeblue}{HTML}{0B5CAD}
\definecolor{codered}{HTML}{B42318}

\setbeamercolor{title}{fg=white,bg=uohblue}
\setbeamercolor{frametitle}{fg=white,bg=uohblue}
\setbeamercolor{block title}{fg=white,bg=uohblue}
\setbeamercolor{block body}{bg=softgray,fg=black}
\setbeamercolor{alerted text}{fg=codered}

\lstdefinestyle{Rstyle}{
  language=R,
  basicstyle=\ttfamily\scriptsize,
  keywordstyle=\color{darkblue}\bfseries,
  commentstyle=\color{codegray}\itshape,
  stringstyle=\color{darkgreen},
  backgroundcolor=\color{codebg},
  frame=single,
  breaklines=true,
  showstringspaces=false,
  columns=fullflexible,
  keepspaces=true,
  upquote=true
}
\lstset{style=Rstyle}


\lstdefinestyle{rcode}{
  language=R,
  basicstyle=\ttfamily\scriptsize,
  keywordstyle=\color{codeblue}\bfseries,
  commentstyle=\color{codegreen},
  stringstyle=\color{codered},
  showstringspaces=false,
  breaklines=true,
  columns=fullflexible,
  keepspaces=true,
  frame=single,
  rulecolor=\color{gray!30},
  backgroundcolor=\color{gray!4}
}
\lstdefinestyle{out}{
  basicstyle=\ttfamily\scriptsize,
  breaklines=true,
  columns=fullflexible,
  keepspaces=true,
  frame=single,
  rulecolor=\color{gray!30},
  backgroundcolor=\color{gray!7}
}


\title[Asociación avanzada]{ Análisis de asociación II}
\subtitle{Filtrado de reglas, visualización e interpretación con R}
\author{Minería de Datos}
\institute{Universidad de O'Higgins}
\date{}

\begin{document}

%------------------------------------------------
\begin{frame}
  \titlepage
\end{frame}

%------------------------------------------------
\begin{frame}{Objetivos de la clase}
\begin{enumerate}
  \item Entender por qué Apriori puede generar miles de reglas y cómo reducirlas.
  \item Filtrar reglas con métricas objetivas, plantillas y redundancia.
  \item Visualizar reglas con \texttt{arulesViz}: scatterplot, matriz, grafo y explorador interactivo.
  \item Interpretar reglas sin confundir la correlación, la causalidad y los artefactos de soporte bajo.
  \item Introducir la discretización, las jerarquías de ítems y el análisis multinivel.
  \item Ejemplos con dataset Zoo
\end{enumerate}
\end{frame}

%------------------------------------------------
\begin{frame}{Flujo de análisis recomendado}
\centering
\begin{tikzpicture}[node distance=0.9cm, every node/.style={font=\small}]
  \node[draw, rounded corners, fill=blue!8, minimum width=2.6cm, minimum height=0.7cm] (a) {1. Codificar datos};
  \node[draw, rounded corners, fill=blue!8, right=of a, minimum width=2.5cm, minimum height=0.7cm] (b) {2. Minar reglas};
  \node[draw, rounded corners, fill=blue!8, right=of b, minimum width=2.9cm, minimum height=0.7cm] (c) {3. Calcular métricas};
  \node[draw, rounded corners, fill=orange!15, below=of b, minimum width=2.9cm, minimum height=0.7cm] (d) {4. Filtrar};
  \node[draw, rounded corners, fill=orange!15, left=of d, minimum width=2.9cm, minimum height=0.7cm] (e) {5. Visualizar};
  \node[draw, rounded corners, fill=green!15, right=of d, minimum width=2.9cm, minimum height=0.7cm] (f) {6. Interpretar};
  \draw[-{Latex}] (a) -- (b);
  \draw[-{Latex}] (b) -- (c);
  \draw[-{Latex}] (c) -- (d);
  \draw[-{Latex}] (d) -- (e);
  \draw[-{Latex}] (d) -- (f);
%  \draw[-{Latex}] (e) -- (f);
\end{tikzpicture}

\vspace{0.4cm}
\textbf{Idea clave:} no se reportan ``todas'' las reglas; se construye un conjunto pequeño, robusto e interpretable.
\end{frame}

%------------------------------------------------
\section{Motivación y métricas}

\begin{frame}{El problema: explosión de reglas}
\begin{block}{Qué ocurre en la práctica}
Un umbral de soporte bajo y confianza alta puede generar miles de reglas, muchas redundantes o demasiado específicas.
\end{block}

\begin{exampleblock}{Ejemplo típico}
\begin{itemize}
  \item \texttt{Groceries}: 9835 transacciones y 169 categorías de producto.
  \item Con soporte bajo aparecen patrones raros, algunos interesantes y otros espurios.
  \item El objetivo es pasar de ``muchas reglas'' a ``pocas reglas accionables''.
\end{itemize}
\end{exampleblock}
\end{frame}

%------------------------------------------------
\begin{frame}{Medidas básicas y avanzadas}
Para una regla $X \Rightarrow Y$:
\begin{align*}
\text{soporte} &= P(X \cap Y) \\
\text{cobertura} &= P(X) \\
\text{confianza} &= P(Y \mid X)=\frac{P(X \cap Y)}{P(X)} \\
\text{lift} &= \frac{P(X \cap Y)}{P(X)P(Y)}=\frac{\text{confianza}}{P(Y)} \\
\text{leverage} &= P(X \cap Y)-P(X)P(Y) \\
\text{conviction} &= \frac{1-P(Y)}{1-\text{confianza}}
\end{align*}
%\end{frame}

%------------------------------------------------
%\begin{frame}{Cómo leer las métricas}
\scriptsize
\begin{tabularx}{\textwidth}{>{\bfseries}lX}
\toprule
Métrica & Interpretación práctica \\
\midrule
Soporte & Prevalencia de la regla. Reglas con soporte muy bajo pueden ser inestables. \\
Confianza & Proporción de transacciones con $X$ que también tienen $Y$. Puede ser engañosa si $Y$ es muy frecuente. \\
Lift & Compara la co-ocurrencia observada contra independencia. Lift $>1$ sugiere asociación positiva. \\
Leverage & Diferencia absoluta entre lo observado y lo esperado por independencia. Útil para priorizar impacto. \\
Conviction & Penaliza contraejemplos de la regla. Crece cuando $X$ rara vez aparece sin $Y$. \\
Count & Número absoluto de transacciones. Siempre revisar junto con soporte. \\
\bottomrule
\end{tabularx}
\end{frame}

%------------------------------------------------
%\begin{frame}[fragile]{Preparación en R}
%\begin{lstlisting}
%# Paquetes principales
%library(arules)
%library(arulesViz)
%library(tidyverse)
%
%# Dataset de mercado incluido en arules
%data("Groceries")
%summary(Groceries)
%itemFrequencyPlot(Groceries, topN = 20, type = "relative")
%\end{lstlisting}
%
%\begin{block}{Nota para R 4.4.x}
%Si CRAN no instala \texttt{arules}, usar una versión archivada compatible:
%\end{block}
%\begin{lstlisting}
%install.packages("remotes", repos = "https://cloud.r-project.org")
%remotes::install_version("arules", version = "1.7-9",
%                         repos = "https://cloud.r-project.org")
%\end{lstlisting}
%\end{frame}

%------------------------------------------------
\section{Filtrado de reglas}

\begin{frame}[fragile]{1. Minar reglas con parámetros iniciales}
\begin{lstlisting}[style=rcode]
rules <- apriori(
  Groceries,
  parameter = list(
    support = 0.005,
    confidence = 0.3,
    minlen = 2,
    maxlen = 5
  )
)

rules
summary(rules)
inspect(head(sort(rules, by = "lift"), 10))
\end{lstlisting}

\begin{itemize}
  \item \texttt{support}: controla prevalencia mínima.
  \item \texttt{confidence}: controla precisión condicional mínima.
  \item \texttt{minlen/maxlen}: evita reglas demasiado simples o demasiado específicas.
\end{itemize}
\end{frame}

%------------------------------------------------
\begin{frame}[fragile]{2. Agregar medidas de interés}
\begin{lstlisting}[style=rcode]
quality(rules) <- cbind(
  quality(rules),
  interestMeasure(
    rules,
    measure = c("leverage", "conviction", "phi", "chiSquared"),
    transactions = Groceries
  )
)

head(quality(rules))
\end{lstlisting}

\begin{block}{Uso práctico}
Ordenar solo por lift puede sobrerrepresentar reglas raras. Combinar lift con soporte, count, leverage o significancia estadística.
\end{block}
\end{frame}

%------------------------------------------------
\begin{frame}[fragile]{3. Filtrado por calidad}
\begin{lstlisting}[style=rcode]
rules_q <- subset(
  rules,
  subset = support >= 0.01 &
           confidence >= 0.5 &
           lift >= 1.5 &
           count >= 50
)

rules_q <- sort(rules_q, by = "lift", decreasing = TRUE)
inspect(head(rules_q, 10))
\end{lstlisting}

\begin{exampleblock}{Criterio defendible}
``Reportamos reglas con soporte $\geq 1\%$, confianza $\geq 50\%$, lift $\geq 1.5$ y al menos 50 transacciones.''
\end{exampleblock}
\end{frame}

%------------------------------------------------
\begin{frame}[fragile]{4. Filtrado por plantillas: RHS específico}
\begin{lstlisting}[style=rcode]
rules_milk <- apriori(
  Groceries,
  parameter = list(support = 0.005, confidence = 0.3, minlen = 2),
  appearance = list(
    rhs = "whole milk",
    default = "lhs"
  )
)

rules_milk <- sort(rules_milk, by = "lift")
inspect(head(rules_milk, 10))
\end{lstlisting}

\begin{block}{Interpretación}
Útil cuando el objetivo es explicar un producto, clase o evento específico: ``qué antecedentes predicen leche entera''. 
\end{block}
\end{frame}

%------------------------------------------------
\begin{frame}[fragile]{5. Filtrado por ítems obligatorios o no deseados}
\begin{lstlisting}[style=rcode]
# Reglas que contienen yogurt en algún lado
rules_yogurt <- subset(rules, items %in% "yogurt")

# Reglas con yogurt en el antecedente y lift alto
rules_yogurt_lhs <- subset(
  rules,
  subset = lhs %in% "yogurt" & lift > 1.5
)

# Reglas que NO contienen soda
rules_no_soda <- subset(rules, !(items %in% "soda"))
\end{lstlisting}

\begin{itemize}
  \item \texttt{items \%in\%}: busca en antecedente o consecuente.
  \item \texttt{lhs \%in\%}: restringe el antecedente.
  \item \texttt{rhs \%in\%}: restringe el consecuente.
\end{itemize}
\end{frame}

%------------------------------------------------
\begin{frame}[fragile]{6. Eliminar reglas redundantes}
\begin{lstlisting}[style=rcode]
length(rules_q)

redundantes <- is.redundant(rules_q)
table(redundantes)

rules_nr <- rules_q[!redundantes]
rules_nr <- sort(rules_nr, by = "lift")

length(rules_nr)
inspect(head(rules_nr, 10))
\end{lstlisting}

\begin{block}{Definición práctica}
Una regla es redundante si existe una regla más general con igual o mayor confianza.
\end{block}
\end{frame}

%------------------------------------------------
\begin{frame}{Ejemplo conceptual de redundancia}
\begin{columns}
\column{0.55\textwidth}
\small
\begin{align*}
\{pan\} &\Rightarrow \{leche\},\quad conf=0.75 \\
\{pan, mantequilla\} &\Rightarrow \{leche\},\quad conf=0.75
\end{align*}

La segunda regla agrega una condición, pero no mejora la confianza.

\column{0.45\textwidth}
\begin{exampleblock}{Conclusión}
Preferir la regla más simple:
\[
\{pan\}\Rightarrow\{leche\}
\]
si su desempeño es igual o mejor.
\end{exampleblock}
\end{columns}
\end{frame}

%------------------------------------------------
\begin{frame}[fragile]{7. Reglas compactas: itemsets cerrados y maximales}
\begin{lstlisting}[style=rcode]
itemsets <- apriori(
  Groceries,
  parameter = list(support = 0.01, target = "frequent itemsets")
)

closed_itemsets <- itemsets[is.closed(itemsets)]
max_itemsets    <- itemsets[is.maximal(itemsets)]
#vector con comparación de items
c(
  frequent = length(itemsets),
  closed   = length(closed_itemsets),
  maximal  = length(max_itemsets)
)
\end{lstlisting}

\begin{itemize}
  \item \textbf{Cerrados}: conservan información de soporte con menos itemsets.
  \item \textbf{Maximales}: resumen aún más, pero pierden soportes de subconjuntos.
\end{itemize}
\end{frame}

%------------------------------------------------
\begin{frame}[fragile]{8. Prueba estadística para una regla}
\begin{lstlisting}[style=rcode]
# Ejemplo: yogurt => whole milk
M <- as(Groceries, "matrix")
X <- M[, "yogurt"]
Y <- M[, "whole milk"]

tab <- table(X, Y)
tab

fisher.test(tab)
chisq.test(tab)
\end{lstlisting}

\begin{block}{Uso práctico}
La prueba ayuda a evaluar si la co-ocurrencia difiere de independencia, pero no reemplaza la interpretación ni corrige automáticamente por múltiples pruebas.
\end{block}
\end{frame}


%------------------------------------------------
\begin{frame}[fragile]{Explicación}
\Large
Queremos entender qué hace este código:

\vspace{0.4cm}
\begin{lstlisting}[style=Rstyle]
# Ejemplo: yogurt => whole milk
M <- as(Groceries, "matrix")
X <- M[, "yogurt"]
Y <- M[, "whole milk"]
tab <- table(X, Y)
tab
fisher.test(tab)
chisq.test(tab)
\end{lstlisting}

\vspace{0.3cm}
\normalsize
La pregunta estadística es:

\begin{center}
\textbf{¿La compra de yogurt está asociada con la compra de leche entera?}
\end{center}
\end{frame}

%------------------------------------------------
\begin{frame}{Regla de asociación evaluada}
\Large
La regla es:

\[
\texttt{yogurt} \Rightarrow \texttt{whole milk}
\]

\vspace{0.3cm}
\normalsize
Interpretación :

\begin{quote}
Si una transacción contiene \texttt{yogurt}, entonces es más probable que también contenga \texttt{whole milk}.
\end{quote}

\vspace{0.3cm}
Pero esta afirmación debe evaluarse con dos criterios:

\begin{enumerate}
  \item \textbf{Magnitud de la asociación}: soporte, confianza, lift.
  \item \textbf{Evidencia estadística}: test de independencia.
\end{enumerate}
\end{frame}

%------------------------------------------------
\begin{frame}[fragile]{Paso 1: convertir transacciones a matriz binaria}
\begin{lstlisting}[style=Rstyle]
M <- as(Groceries, "matrix")
\end{lstlisting}

\vspace{0.3cm}
\texttt{Groceries} es un conjunto de transacciones. Cada fila es una compra y cada columna es un producto.

\vspace{0.3cm}
Después de convertirlo a matriz:

\[
M_{ij} =
\begin{cases}
\texttt{TRUE}, & \text{si la transacción } i \text{ contiene el producto } j \\
\texttt{FALSE}, & \text{si no lo contiene}
\end{cases}
\]

\vspace{0.3cm}
Ejemplo conceptual:

\begin{center}
\begin{tabular}{lccc}
\toprule
Transacción & yogurt & whole milk & bread \\
\midrule
T1 & TRUE  & TRUE  & FALSE \\
T2 & FALSE & TRUE  & TRUE \\
T3 & TRUE  & FALSE & TRUE \\
\bottomrule
\end{tabular}
\end{center}
\end{frame}

%------------------------------------------------
\begin{frame}[fragile]{Paso 2: definir antecedente y consecuente}
\begin{lstlisting}[style=Rstyle]
X <- M[, "yogurt"]
Y <- M[, "whole milk"]
\end{lstlisting}

\vspace{0.4cm}
\Large
\[
X = \texttt{yogurt}
\]

\[
Y = \texttt{whole milk}
\]

\vspace{0.3cm}
\normalsize
Cada variable es un vector lógico:

\begin{itemize}
  \item \texttt{X = TRUE}: la transacción contiene yogurt.
  \item \texttt{Y = TRUE}: la transacción contiene leche entera.
\end{itemize}

\vspace{0.3cm}
Por lo tanto, el código transforma una regla en un problema clásico de asociación entre dos variables binarias.
\end{frame}

%------------------------------------------------
\begin{frame}[fragile]{Paso 3: construir la tabla de contingencia}
\begin{lstlisting}[style=Rstyle]
tab <- table(X, Y)
tab
\end{lstlisting}

\vspace{0.2cm}
Salida  para \texttt{Groceries}:

\begin{center}
\begin{tabular}{lrr}
\toprule
 & \texttt{Y = FALSE} & \texttt{Y = TRUE} \\
\midrule
\texttt{X = FALSE} & 6501 & 1962 \\
\texttt{X = TRUE}  & 821  & 551 \\
\bottomrule
\end{tabular}
\end{center}

\vspace{0.2cm}
Donde:

\begin{itemize}
  \item 551 compras contienen \texttt{yogurt} y \texttt{whole milk}.
  \item 821 contienen \texttt{yogurt}, pero no \texttt{whole milk}.
  \item 1962 contienen \texttt{whole milk}, pero no \texttt{yogurt}.
  \item 6501 no contienen ninguno de los dos productos.
\end{itemize}
\end{frame}

%------------------------------------------------
\begin{frame}{Tabla observada: notación estadística}
\Large
\[
\begin{array}{c|cc|c}
 & Y=0 & Y=1 & \text{Total} \\
\hline
X=0 & d=6501 & c=1962 & 8463 \\
X=1 & b=821  & a=551  & 1372 \\
\hline
\text{Total} & 7322 & 2513 & 9835
\end{array}
\]

\vspace{0.3cm}
\normalsize
Para la regla:

\[
X \Rightarrow Y
\]

el conteo más importante es:

\[
a = \#(X=1, Y=1) = 551
\]

porque corresponde a las transacciones que contienen ambos productos.
\end{frame}

%------------------------------------------------
\begin{frame}{Métricas de la regla}
Con los conteos anteriores:

\[
N = 9835, \quad a=551, \quad b=821, \quad c=1962, \quad d=6501
\]

\vspace{0.2cm}
\begin{align*}
\text{soporte}(X \Rightarrow Y) &= P(X,Y) = \frac{551}{9835} = 0.056 \\
\text{confianza}(X \Rightarrow Y) &= P(Y\mid X) = \frac{551}{551+821} = 0.402 \\
\text{soporte}(Y) &= P(Y) = \frac{2513}{9835} = 0.256 \\
\text{lift}(X \Rightarrow Y) &= \frac{P(Y\mid X)}{P(Y)} = \frac{0.402}{0.256} = 1.57
\end{align*}

\vspace{0.2cm}
\textbf{Interpretación:} entre quienes compran yogurt, la probabilidad de comprar leche entera es 1.57 veces la probabilidad basal de comprar leche entera.
\end{frame}

%------------------------------------------------
\begin{frame}{¿Por qué hacer un test estadístico?}
\Large
Las métricas describen la asociación, pero no responden completamente:

\vspace{0.4cm}
\begin{center}
\textbf{¿La co-ocurrencia observada podría explicarse por azar?}
\end{center}

\vspace{0.4cm}
\normalsize
Para eso se testea independencia:

\[
H_0: X \text{ e } Y \text{ son independientes}
\]

\[
H_1: X \text{ e } Y \text{ no son independientes}
\]

\vspace{0.3cm}
En términos de la regla:

\[
H_0: P(Y\mid X) = P(Y)
\]

\[
H_1: P(Y\mid X) \neq P(Y)
\]
\end{frame}

%------------------------------------------------
\begin{frame}{Conteos esperados bajo independencia}
Si \texttt{yogurt} y \texttt{whole milk} fueran independientes, el conteo esperado en cada celda sería:

\[
E_{ij} = \frac{(\text{total fila } i)(\text{total columna } j)}{N}
\]

\vspace{0.2cm}
\begin{center}
\begin{tabular}{lrr}
\toprule
 & \texttt{Y = FALSE} & \texttt{Y = TRUE} \\
\midrule
\texttt{X = FALSE} & 6300.57 & 2162.43 \\
\texttt{X = TRUE}  & 1021.43 & 350.57 \\
\bottomrule
\end{tabular}
\end{center}

\vspace{0.3cm}
Celda clave:

\[
E(X=1,Y=1)=350.57
\]

Pero observamos:

\[
O(X=1,Y=1)=551
\]

Hay muchas más compras conjuntas de \texttt{yogurt} y \texttt{whole milk} de las esperadas bajo independencia.
\end{frame}

%------------------------------------------------
\begin{frame}[fragile]{Test chi-cuadrado de independencia}
\begin{lstlisting}[style=Rstyle]
chisq.test(tab)
\end{lstlisting}

\vspace{0.3cm}
El test compara conteos observados contra conteos esperados:

\[
\chi^2 = \sum \frac{(O_{ij} - E_{ij})^2}{E_{ij}}
\]


\begin{lstlisting}[style=Rstyle]
Pearson's Chi-squared test with Yates' continuity correction

X-squared = 177.99, df = 1, p-value < 2.2e-16
\end{lstlisting}

\vspace{0.2cm}
\textbf{Interpretación:} el valor-p es extremadamente pequeño, por lo tanto se rechaza la independencia entre \texttt{yogurt} y \texttt{whole milk}.


\end{frame}

%------------------------------------------------
\begin{frame}[fragile]{Test exacto de Fisher}
\begin{lstlisting}[style=Rstyle]
fisher.test(tab)
\end{lstlisting}

\vspace{0.3cm}
El test exacto de Fisher también evalúa independencia en una tabla $2\times2$.

\vspace{0.2cm}
Es especialmente útil cuando hay conteos pequeños. En este ejemplo los conteos son grandes, pero Fisher sigue siendo válido.


\begin{lstlisting}[style=Rstyle]
Fisher's Exact Test for Count Data

p-value < 2.2e-16
alternative hypothesis: true odds ratio is not equal to 1
sample estimates:
odds ratio = 2.22
\end{lstlisting}

\vspace{0.2cm}
\textbf{Interpretación:} las odds de comprar \texttt{whole milk} son mayores entre quienes compran \texttt{yogurt}.
\end{frame}

%------------------------------------------------
\begin{frame}{Odds ratio y asociación positiva}
El odds ratio de la tabla es:

\[
OR = \frac{a \cdot d}{b \cdot c}
\]

\[
OR = \frac{551 \times 6501}{821 \times 1962} = 2.22
\]

\vspace{0.4cm}
Interpretación:

\begin{itemize}
  \item $OR = 1$: no hay asociación.
  \item $OR > 1$: asociación positiva.
  \item $OR < 1$: asociación negativa.
\end{itemize}

\vspace{0.3cm}
En este caso:

\[
OR = 2.22 > 1
\]

lo que apoya una asociación positiva entre \texttt{yogurt} y \texttt{whole milk}.

\begin{quote}
La regla \texttt{yogurt} $\Rightarrow$ \texttt{whole milk} muestra una asociación positiva y estadísticamente significativa en el dataset \texttt{Groceries}.
\end{quote}

\end{frame}

%------------------------------------------------
\begin{frame}[fragile]{Test direccional: asociación positiva}
Por defecto, Fisher testea si existe asociación en cualquier dirección:

\[
H_1: OR \neq 1
\]

\vspace{0.3cm}
Pero si queremos testear específicamente una regla positiva:

\[
\texttt{yogurt} \Rightarrow \texttt{whole milk}
\]

podemos usar:

\begin{lstlisting}[style=Rstyle]
fisher.test(tab, alternative = "greater")
\end{lstlisting}

\vspace{0.2cm}
Esto evalúa:

\[
H_1: OR > 1
\]

\vspace{0.2cm}
Es decir, si \texttt{yogurt} y \texttt{whole milk} aparecen juntos más de lo esperado bajo independencia.
\end{frame}

%------------------------------------------------
\begin{frame}{Cuidado con la interpretación}
\Large
El test demuestra asociación, no causalidad.

\vspace{0.4cm}
\normalsize
No podemos concluir:

\begin{quote}
Comprar yogurt causa comprar leche entera.
\end{quote}

\vspace{0.3cm}
Sí podemos concluir:

\begin{quote}
En estas transacciones, yogurt y leche entera aparecen juntos más de lo esperado bajo independencia.
\end{quote}

\vspace{0.4cm}
Además, en minería de reglas se prueban muchas asociaciones; por lo tanto, en análisis formales conviene considerar ajuste por múltiples pruebas o usar el test como herramienta interpretativa secundaria.
\end{frame}







%------------------------------------------------
\begin{frame}[fragile]{9. Barrido de parámetros}
\begin{lstlisting}[style=rcode]
grid <- expand.grid(
  support = c(0.001, 0.005, 0.01, 0.02),
  confidence = c(0.3, 0.5, 0.7)
)

grid$n_rules <- purrr::map2_int(
  grid$support, grid$confidence,
  ~ length(apriori(
      Groceries,
      parameter = list(support = .x, confidence = .y, minlen = 2),
      control = list(verbose = FALSE)
    ))
)

grid
\end{lstlisting}
\end{frame}

%------------------------------------------------
\begin{frame}[fragile]{Visualizar sensibilidad de parámetros}
\begin{lstlisting}[style=rcode]
ggplot(grid, aes(x = support, y = n_rules, color = factor(confidence))) +
  geom_line() +
  geom_point() +
  scale_x_log10() +
  labs(
    x = "Soporte minimo",
    y = "Numero de reglas",
    color = "Confianza"
  ) +
  theme_minimal()
\end{lstlisting}

\begin{block}{Lectura}
Una caída abrupta indica que el umbral de soporte está controlando fuertemente el espacio de búsqueda.
\end{block}
\end{frame}


%------------------------------------------------
\section{Visualizacion}
\begin{frame}[fragile]{Visualizacion 1: frecuencia de items}
\begin{columns}[T]
\column{0.47\textwidth}
\begin{lstlisting}[style=rcode]
itemFrequencyPlot(trans, topN = 20)
\end{lstlisting}

\begin{block}{Que mirar}
\begin{itemize}
  \item Items muy frecuentes: potenciales reglas triviales.
  \item Items raros: requieren soporte bajo.
  \item Variables de clase: \texttt{type=...}.
\end{itemize}
\end{block}

\column{0.49\textwidth}
\begin{tikzpicture}
\begin{axis}[
  xbar,
  width=0.98\textwidth,
  height=0.58\textheight,
  xmin=0, xmax=90,
  xlabel={count},
  symbolic y coords={toothed,tail,legs,breathes,backbone},
  ytick=data,
  nodes near coords,
  nodes near coords align={horizontal},
]
\addplot coordinates {(61,toothed) (75,tail) (78,legs) (80,breathes) (83,backbone)};
\end{axis}
\end{tikzpicture}
\end{columns}
\end{frame}

%------------------------------------------------
\begin{frame}[fragile]{Visualizacion 2: scatterplot de reglas}
\begin{columns}[T]
\column{0.46\textwidth}
\begin{lstlisting}[style=rcode]
plot(rules)
plot(rules, shading = "order")
plot(rules, control = list(jitter = 0))
\end{lstlisting}

\begin{block}{Como leerlo}
\begin{itemize}
  \item X: soporte.
  \item Y: confianza.
  \item Color: lift u orden de la regla.
  \item Nubes densas indican reglas repetitivas.
\end{itemize}
\end{block}

\column{0.50\textwidth}
\begin{tikzpicture}
\begin{axis}[
  width=0.98\textwidth,
  height=0.58\textheight,
  xlabel={support}, ylabel={confidence},
  xmin=0.04, xmax=0.45, ymin=0.75, ymax=1.02,
  grid=both,
]
\addplot[only marks, mark=*, mark size=1.8pt] coordinates {
(0.079,1.00) (0.129,1.00) (0.198,1.00) (0.406,1.00)
(0.386,0.907) (0.158,0.889) (0.129,0.813) (0.297,1.00)
(0.327,1.00) (0.376,1.00) (0.168,0.850) (0.139,0.824)
};
\end{axis}
\end{tikzpicture}
\end{columns}
\end{frame}

%------------------------------------------------
\begin{frame}[fragile]{Visualizacion 3: matriz agrupada}
\begin{columns}[T]
\column{0.47\textwidth}
\begin{lstlisting}[style=rcode]
set.seed(1234)
plot(rules,
     method = "grouped matrix")
\end{lstlisting}

\begin{block}{Como leerla}
Agrupa reglas con relaciones LHS-RHS similares. Es util cuando hay muchas reglas, porque revela bloques: peces, aves, mamiferos.
\end{block}

\column{0.49\textwidth}
\scriptsize
\begin{tabular}{lccc}
\toprule
Grupo de reglas & RHS dominante & Rasgos frecuentes & Interpretacion \\
\midrule
G1 & type=fish & eggs, fins, aquatic & Perfil acuatico \\
G2 & type=bird & feathers, airborne, tail & Perfil ave \\
G3 & type=mammal & milk, hair, toothed & Perfil mamifero \\
\bottomrule
\end{tabular}

\vspace{0.35cm}
\begin{center}
\begin{tikzpicture}[scale=0.85]
\foreach \x/\y/\s/\c in {0/0/1.0/blue,1/0/0.75/blue,2/0/0.5/blue,0/1/0.25/gray,1/1/0.9/orange,2/1/0.6/orange,0/2/0.3/gray,1/2/0.5/gray,2/2/0.95/green} {
  \fill[\c!55] (\x,\y) circle (0.25*\s+0.05);
}
\node at (0,-0.55) {fish}; \node at (1,-0.55) {bird}; \node at (2,-0.55) {mammal};
\node[rotate=90] at (-0.55,0) {fins}; \node[rotate=90] at (-0.55,1) {airborne}; \node[rotate=90] at (-0.55,2) {milk};
\end{tikzpicture}
\end{center}
\end{columns}
\end{frame}

%------------------------------------------------
\begin{frame}[fragile]{Visualizacion 4: grafo de reglas}
\begin{columns}[T]
\column{0.46\textwidth}
\begin{lstlisting}[style=rcode]
plot(rules, method = "graph")

plot(head(rules, by = "phi", n = 100),
     method = "graph")
\end{lstlisting}

\begin{block}{Cuando usarlo}
Usar pocos nodos: top 50 o top 100 por lift, phi o confianza. Si se grafican miles de reglas, el grafo deja de ser interpretable.
\end{block}

\column{0.50\textwidth}
\begin{center}
\begin{tikzpicture}[node distance=1.0cm, every node/.style={font=\scriptsize}]
\node[draw, circle, fill=green!15] (milk) {milk};
\node[draw, circle, fill=green!15, below left=0.75cm of milk] (hair) {hair};
\node[draw, rounded corners, fill=green!30, right=1.2cm of milk] (mammal) {type=mammal};
\node[draw, circle, fill=orange!15, above=0.75cm of mammal] (feath) {feathers};
\node[draw, rounded corners, fill=orange!30, right=1.3cm of feath] (bird) {type=bird};
\node[draw, circle, fill=blue!15, below=0.75cm of mammal] (fins) {fins};
\node[draw, circle, fill=blue!15, below=0.75cm of fins] (eggs) {eggs};
\node[draw, rounded corners, fill=blue!30, right=1.3cm of fins] (fish) {type=fish};
\draw[->, thick] (milk) -- node[above]{1.0} (mammal);
\draw[->, thick] (hair) -- (mammal);
\draw[->, thick] (feath) -- node[above]{1.0} (bird);
\draw[->, thick] (fins) -- (fish);
\draw[->, thick] (eggs) -- (fish);
\end{tikzpicture}
\end{center}
\end{columns}
\end{frame}

%%------------------------------------------------

\section{Interpretación}

\begin{frame}{Errores frecuentes de interpretación}
\begin{enumerate}
  \item \textbf{Confundir confianza con causalidad}: $X \Rightarrow Y$ no significa que $X$ cause $Y$.
  \item \textbf{Ignorar la frecuencia base de $Y$}: si $Y$ es muy común, la confianza puede ser alta sin asociación real.
  \item \textbf{Sobrevalorar lift con soporte bajo}: una regla rara puede tener lift extremo por azar.
  \item \textbf{No revisar conteos absolutos}: soporte relativo sin count puede ocultar inestabilidad.
  \item \textbf{Reportar reglas redundantes}: reglas más largas no siempre son mejores.
\end{enumerate}
\end{frame}

%------------------------------------------------
\begin{frame}{Recomendaciones para interpretar una regla}
\begin{exampleblock}{Regla}
\[
\{yogurt, other\ vegetables\} \Rightarrow \{whole\ milk\}
\]
\end{exampleblock}

\begin{enumerate}
  \item \textbf{Cobertura}: qué proporción compra yogurt y verduras.
  \item \textbf{Confianza}: entre quienes compran eso, cuántos también compran leche.
  \item \textbf{Lift}: cuánto aumenta la probabilidad respecto a comprar leche al azar.
  \item \textbf{Count}: cuántas transacciones reales apoyan la regla.
  \item \textbf{Acción}: promoción cruzada, layout, recomendador o hipótesis de comportamiento.
\end{enumerate}
\end{frame}

%%------------------------------------------------
%\begin{frame}{Decisión: ¿reportar o descartar?}
%\scriptsize
%\begin{tabularx}{\textwidth}{lXXX}
%\toprule
%Caso & Patrón & Riesgo & Decisión sugerida \\
%\midrule
%Alta confianza, lift $\approx 1$ & $Y$ es muy frecuente & Regla trivial & Descartar o usar como baseline \\
%Lift alto, count bajo & Asociación rara & Azar / inestabilidad & Requiere validación \\
%Soporte alto, lift moderado & Patrón común & Puede ser accionable & Reportar si interpreta negocio \\
%Regla larga & Muy específica & Redundancia / sobreajuste & Comparar con reglas más generales \\
%RHS objetivo & Regla explicativa & Sesgo de selección & Declarar plantilla usada \\
%\bottomrule
%\end{tabularx}
%\end{frame}
%
%%------------------------------------------------
%\begin{frame}[fragile]{Exportar reglas para informe}
%\begin{lstlisting}
%rules_tbl <- as(rules_nr, "data.frame") %>%
%  arrange(desc(lift)) %>%
%  select(rules, support, confidence, coverage, lift, count, leverage)
%
%head(rules_tbl, 20)
%
%write.csv(rules_tbl, "reglas_filtradas.csv", row.names = FALSE)
%\end{lstlisting}
%
%\begin{block}{Para reportes}
%Incluir reglas, métricas, criterio de filtrado y una breve interpretación por grupo de reglas.
%\end{block}
%\end{frame}

%------------------------------------------------
\section{Datos no transaccionales y jerarquías}

\begin{frame}{Variables categóricas, lógicas y continuas}
\begin{itemize}
  \item Variables categóricas: cada nivel se transforma en un ítem, por ejemplo \texttt{Species=setosa}.
  \item Variables lógicas: normalmente solo \texttt{TRUE} genera un ítem.
  \item Variables continuas: deben discretizarse antes de transformarse en ítems.
\end{itemize}

\begin{block}{Riesgo metodológico}
La discretización cambia las reglas encontradas. Debe justificarse y, si es posible, evaluarse con sensibilidad.
\end{block}
\end{frame}

%------------------------------------------------
\begin{frame}[fragile]{Ejemplo: discretización con iris}
\begin{lstlisting}[style=rcode]
data(iris)
iris$Versicolor <- iris$Species == "versicolor"

iris_disc <- iris %>%
  mutate(
    Petal.Length = discretize(Petal.Length,
                              method = "frequency",
                              breaks = 3,
                              labels = c("short", "medium", "long")),
    Petal.Width = discretize(Petal.Width,
                             method = "frequency",
                             breaks = 2,
                             labels = c("narrow", "wide"))
  )

iris_trans <- transactions(iris_disc)
itemLabels(iris_trans)
\end{lstlisting}
\end{frame}

%------------------------------------------------
\begin{frame}[fragile]{Reglas de asociación orientadas a clase}
\begin{lstlisting}[style=rcode]
iris_rules <- apriori(
  iris_trans,
  parameter = list(support = 0.05, confidence = 0.8, minlen = 2),
  appearance = list(
    rhs = c("Species=setosa", "Species=versicolor", "Species=virginica"),
    default = "lhs"
  )
)

iris_rules <- iris_rules[!is.redundant(iris_rules)]
inspect(sort(iris_rules, by = "lift"))
\end{lstlisting}

\begin{block}{Uso}
Puede verse como un clasificador interpretable basado en reglas.
\end{block}
\end{frame}

%------------------------------------------------
\begin{frame}{Jerarquías de ítems}
\begin{exampleblock}{Ejemplo supermercado}
\begin{itemize}
  \item \texttt{whole milk}, \texttt{yogurt}, \texttt{cream} $\rightarrow$ \texttt{dairy produce}
  \item \texttt{citrus fruit}, \texttt{tropical fruit} $\rightarrow$ \texttt{fruit}
\end{itemize}
\end{exampleblock}

\begin{block}{Ventaja}
Permite descubrir reglas a nivel de producto, categoría o ambos niveles simultáneamente.
\end{block}
\end{frame}

%------------------------------------------------
\begin{frame}[fragile]{Agregación de ítems en Groceries}
\begin{lstlisting}[style=rcode]
data("Groceries")
head(itemInfo(Groceries))

Groceries_level2 <- aggregate(Groceries, by = "level2")
Groceries_level2

inspect(head(Groceries, 3))
inspect(head(Groceries_level2, 3))

rules_level2 <- apriori(
  Groceries_level2,
  parameter = list(support = 0.005, confidence = 0.8)
)
inspect(head(sort(rules_level2, by = "lift"), 10))
\end{lstlisting}
\end{frame}

%------------------------------------------------
\begin{frame}[fragile]{Análisis multinivel}
\begin{lstlisting}[style=rcode]
Groceries_multi <- addAggregate(Groceries, "level2")
inspect(head(Groceries_multi, 3))

rules_multi <- apriori(
  Groceries_multi,
  parameter = list(support = 0.005, confidence = 0.8)
)

# Elimina reglas espurias: item A => grupo de item A
rules_multi <- filterAggregate(rules_multi)
inspect(head(sort(rules_multi, by = "lift"), 10))
\end{lstlisting}

\begin{block}{Interpretación}
Ayuda a estudiar si la asociación ocurre entre productos concretos o entre categorías generales.
\end{block}
\end{frame}



%------------------------------------------------
\begin{frame}[fragile]{Dataset Zoo: animales como transacciones}
\begin{columns}[T]
\column{0.52\textwidth}
\begin{lstlisting}[style=rcode]
library(arules)
library(arulesViz)
library(mlbench)
library(tidyverse)

data(Zoo, package = "mlbench")

Zoo_has_legs <- Zoo |>
  mutate(legs = legs > 0)

trans <- transactions(Zoo_has_legs)
trans
\end{lstlisting}

\column{0.44\textwidth}
\begin{lstlisting}[style=out]
transactions in sparse format with
 101 transactions (rows)
 23 items (columns)
\end{lstlisting}

\begin{block}{Interpretacion}
Cada animal es una transaccion. Cada rasgo es un item: \texttt{hair}, \texttt{milk}, \texttt{fins}, \texttt{type=mammal}, etc.
\end{block}
\end{columns}
\end{frame}

%------------------------------------------------
\begin{frame}[fragile]{Antes de minar reglas: inspeccionar las transacciones}
\begin{columns}[T]
\column{0.54\textwidth}
\begin{lstlisting}[style=rcode]
summary(trans)
itemFrequencyPlot(trans, topN = 20)
inspect(trans[1:3])
\end{lstlisting}

\column{0.42\textwidth}
\begin{lstlisting}[style=out]
most frequent items:
backbone breathes legs tail toothed
      83      80   78   75      61

itemset length:
Min. 3 | Median 9 | Mean 8.31
Max. 12
\end{lstlisting}
\end{columns}

\begin{alertblock}{Recordar}
Si los items frecuentes son muy generales, muchas reglas tendran alta confianza pero baja utilidad interpretativa.
\end{alertblock}
\end{frame}

%------------------------------------------------
\begin{frame}[fragile]{El problema: demasiadas reglas}
\begin{columns}[T]
\column{0.50\textwidth}
\begin{lstlisting}[style=rcode]
rules <- apriori(
  trans,
  parameter = list(
    support = 0.05,
    confidence = 0.9
  )
)
length(rules)
\end{lstlisting}

\column{0.46\textwidth}
\begin{lstlisting}[style=out]
Absolute minimum support count: 5
set transactions ...
  [23 item(s), 101 transaction(s)]
checking subsets of size 1 2 ... 10
writing ... [7174 rule(s)]

[1] 7174
\end{lstlisting}
\end{columns}

\vspace{0.2cm}
\begin{block}{Pregunta }
\centering
¿Cómo pasamos de 7174 reglas a un conjunto que se pueda explicar en clase o en un informe?
\end{block}
\end{frame}

%------------------------------------------------
\begin{frame}[fragile]{Mapa general de mecanismos de filtrado}
\begin{center}
\begin{tikzpicture}[node distance=1.2cm, every node/.style={font=\small}]
\node[draw, rounded corners, fill=blue!8, minimum width=3.2cm, minimum height=0.75cm] (raw) {Reglas crudas};
\node[draw, rounded corners, fill=green!10, right=0.8cm of raw, minimum width=3.6cm, minimum height=0.75cm] (metrics) {Metricas: soporte, confianza, lift};
\node[draw, rounded corners, fill=orange!12, right=0.8cm of metrics, minimum width=3.4cm, minimum height=0.75cm] (templates) {Plantillas: RHS/LHS};
\node[draw, rounded corners, fill=purple!10, below=0.8cm of metrics, minimum width=3.5cm, minimum height=0.75cm] (redund) {No redundantes};
\node[draw, rounded corners, fill=red!8, right=0.8cm of redund, minimum width=3.2cm, minimum height=0.75cm] (viz) {Visualizacion};
\node[draw, rounded corners, fill=gray!15, below=0.8cm of redund, minimum width=4.0cm, minimum height=0.75cm] (final) {Reglas interpretables};
\draw[->, thick] (raw) -- (metrics);
\draw[->, thick] (metrics) -- (templates);
\draw[->, thick] (templates) |- (redund);
\draw[->, thick] (redund) -- (viz);
\draw[->, thick] (viz) |- (final);
\end{tikzpicture}
\end{center}

\begin{block}{Idea central}
El filtrado debe combinar criterios estadísticos, criterios computacionales y criterios de interpretación del problema.
\end{block}
\end{frame}

%------------------------------------------------
\section{Filtrado por parámetros}
\begin{frame}[fragile]{Filtro 1: soporte minimo}
\begin{block}{Definicion}
\[
\mathrm{supp}(X \Rightarrow Y) = \frac{\#\{t: X \cup Y \subseteq t\}}{N}
\]
\end{block}

\begin{columns}[T]
\column{0.52\textwidth}
\begin{lstlisting}[style=rcode]
# soporte para al menos 5 animales
5 / nrow(trans)

its_10 <- apriori(trans,
  parameter = list(target = "frequent",
                   support = 0.10))

its_05 <- apriori(trans,
  parameter = list(target = "frequent",
                   support = 0.05))
\end{lstlisting}

\column{0.44\textwidth}
\begin{lstlisting}[style=out]
[1] 0.0495

support = 0.10 -> 1465 itemsets
support = 0.05 -> 2537 itemsets
\end{lstlisting}
\end{columns}

\begin{alertblock}{Interpretacion}
Bajar soporte de 10\% a 5\% permite detectar patrones que aparecen en al menos 5 animales, pero aumenta mucho el espacio de búsqueda.
\end{alertblock}
\end{frame}

%------------------------------------------------
\begin{frame}[fragile]{Resultado del soporte: itemsets frecuentes principales}
\begin{columns}[T]
\column{0.50\textwidth}
\begin{lstlisting}[style=rcode]
its <- sort(its_05, by = "support")
inspect(head(its, n = 10))
\end{lstlisting}

\column{0.46\textwidth}
\scriptsize
\begin{tabular}{lrr}
\toprule
Itemset & support & count \\
\midrule
\{backbone\} & 0.8218 & 83 \\
\{breathes\} & 0.7921 & 80 \\
\{legs\} & 0.7723 & 78 \\
\{tail\} & 0.7426 & 75 \\
\{backbone, tail\} & 0.7327 & 74 \\
\{breathes, legs\} & 0.7228 & 73 \\
\{backbone, breathes\} & 0.6832 & 69 \\
\{backbone, legs\} & 0.6337 & 64 \\
\{backbone, breathes, legs\} & 0.6337 & 64 \\
\{toothed\} & 0.6040 & 61 \\
\bottomrule
\end{tabular}
\end{columns}

\begin{block}{Lectura}
Los itemsets mas frecuentes describen rasgos generales. Son utiles para explorar estructura, pero no necesariamente son las reglas mas informativas.
\end{block}
\end{frame}

%------------------------------------------------
\begin{frame}[fragile]{Filtro 2: confianza minima}
\begin{block}{Definición}
\[
\mathrm{conf}(X \Rightarrow Y) = \frac{\mathrm{supp}(X \cup Y)}{\mathrm{supp}(X)} = P(Y \mid X)
\]
\end{block}

\begin{columns}[T]
\column{0.48\textwidth}
\begin{lstlisting}[style=rcode]
rules <- apriori(trans,
  parameter = list(
    support = 0.05,
    confidence = 0.9))

inspect(head(rules))
\end{lstlisting}

\column{0.48\textwidth}
\scriptsize
\begin{tabular}{llrrr}
\toprule
LHS & RHS & supp. & conf. & lift \\
\midrule
\{type=insect\} & \{eggs\} & 0.079 & 1.0 & 1.712 \\
\{type=insect\} & \{legs\} & 0.079 & 1.0 & 1.295 \\
\{type=insect\} & \{breathes\} & 0.079 & 1.0 & 1.262 \\
\{type=mollusc...\} & \{eggs\} & 0.089 & 0.9 & 1.541 \\
\{type=fish\} & \{fins\} & 0.129 & 1.0 & 5.941 \\
\{type=fish\} & \{aquatic\} & 0.129 & 1.0 & 2.806 \\
\bottomrule
\end{tabular}
\end{columns}

\begin{alertblock}{Cuidado}
Alta confianza no implica relevancia: puede deberse a que el RHS es muy frecuente.
\end{alertblock}
\end{frame}

%------------------------------------------------
\begin{frame}[fragile]{Filtro 3: lift para evitar reglas triviales}
\begin{block}{Definición}
\[
\mathrm{lift}(X \Rightarrow Y) = \frac{\mathrm{conf}(X \Rightarrow Y)}{\mathrm{supp}(Y)}
\]
\end{block}

\begin{columns}[T]
\column{0.45\textwidth}
\begin{lstlisting}[style=rcode]
rules_lift <- sort(rules, by = "lift")
inspect(head(rules_lift, n = 6))
\end{lstlisting}

\column{0.51\textwidth}
\scriptsize
\begin{tabular}{llrrrr}
\toprule
LHS & RHS & supp. & conf. & lift & count \\
\midrule
\{eggs, fins\} & \{type=fish\} & 0.1287 & 1 & 7.769 & 13 \\
\{eggs, aquatic, fins\} & \{type=fish\} & 0.1287 & 1 & 7.769 & 13 \\
\{eggs, toothed, fins\} & \{type=fish\} & 0.1287 & 1 & 7.769 & 13 \\
\{eggs, fins, tail\} & \{type=fish\} & 0.1287 & 1 & 7.769 & 13 \\
\{eggs, backbone, fins\} & \{type=fish\} & 0.1287 & 1 & 7.769 & 13 \\
\{eggs, aquatic, toothed, fins\} & \{type=fish\} & 0.1287 & 1 & 7.769 & 13 \\
\bottomrule
\end{tabular}
\end{columns}

\begin{block}{Interpretacion biologica simple}
En \texttt{Zoo}, los peces tienen un perfil muy distintivo: \texttt{fins}, \texttt{eggs} y rasgos acuaticos predicen fuertemente \texttt{type=fish}.
\end{block}
\end{frame}

%%------------------------------------------------
%\begin{frame}[fragile]{Resumen: que significa cada metrica?}
%\scriptsize
%\begin{tabular}{p{0.20\textwidth}p{0.30\textwidth}p{0.38\textwidth}}
%\toprule
%Metrica & Pregunta que responde & Uso para filtrar \\
%\midrule
%Soporte & Cuantas transacciones contienen la regla? & Elimina patrones demasiado raros o inestables. \\
%Confianza & Cuando ocurre el antecedente, cuantas veces ocurre el consecuente? & Busca reglas predictivas, pero puede favorecer RHS frecuentes. \\
%Coverage & Cuantas transacciones cubre el antecedente? & Evita reglas con antecedente demasiado especifico. \\
%Lift & La regla mejora sobre el azar o frecuencia basal del RHS? & Prioriza reglas no triviales. Lift $>1$ indica asociacion positiva. \\
%Count & Numero absoluto de casos que sustentan la regla. & Facilita discutir evidencia con estudiantes. \\
%\bottomrule
%\end{tabular}
%
%\vspace{0.4cm}
%\begin{alertblock}{Regla pedagogica}
%Para ensenar e interpretar, nunca mostrar solo confianza. Mostrar al menos: soporte, confianza, lift y count.
%\end{alertblock}
%\end{frame}

%------------------------------------------------
\section{Filtrado con plantillas}
\begin{frame}[fragile]{Filtro 4: restringir el RHS con \texttt{appearance}}
\begin{block}{Motivacion}
Si el objetivo es interpretar rasgos que predicen el tipo de animal, conviene forzar reglas de la forma:
\[
\mathrm{rasgos} \Rightarrow \mathrm{type=animal}
\]
\end{block}

\begin{columns}[T]
\column{0.52\textwidth}
\begin{lstlisting}[style=rcode]
type <- grep("type=", itemLabels(trans),
             value = TRUE)

type

rules_type <- apriori(
  trans,
  appearance = list(rhs = type)
)
rules_type
\end{lstlisting}

\column{0.44\textwidth}
\begin{lstlisting}[style=out]
[1] "type=mammal" "type=bird"
[3] "type=reptile" "type=fish"
[5] "type=amphibian" "type=insect"
[7] "type=mollusc.et.al"

writing ... [571 rule(s)]

set of 571 rules
\end{lstlisting}
\end{columns}
\end{frame}

%------------------------------------------------
\begin{frame}[fragile]{Resultado de \texttt{appearance}: reglas orientadas a clase}
\begin{lstlisting}[style=rcode]
rules_type |>
  sort(by = "lift") |>
  head() |>
  inspect()
\end{lstlisting}

\vspace{0.1cm}
\scriptsize
\begin{tabular}{llrrrr}
\toprule
LHS & RHS & support & confidence & lift & count \\
\midrule
\{eggs, fins\} & \{type=fish\} & 0.1287 & 1 & 7.769 & 13 \\
\{eggs, aquatic, fins\} & \{type=fish\} & 0.1287 & 1 & 7.769 & 13 \\
\{eggs, toothed, fins\} & \{type=fish\} & 0.1287 & 1 & 7.769 & 13 \\
\{eggs, fins, tail\} & \{type=fish\} & 0.1287 & 1 & 7.769 & 13 \\
\{eggs, backbone, fins\} & \{type=fish\} & 0.1287 & 1 & 7.769 & 13 \\
\{eggs, aquatic, toothed, fins\} & \{type=fish\} & 0.1287 & 1 & 7.769 & 13 \\
\bottomrule
\end{tabular}

\begin{alertblock}{Interpretación}
El filtro por plantilla elimina reglas como \texttt{type=fish => fins} y enfoca la pregunta en reglas tipo clasificador: rasgos $\Rightarrow$ clase.
\end{alertblock}
\end{frame}

%------------------------------------------------
\begin{frame}[fragile]{Filtro 5: tamaño de la regla}
\begin{block}{Motivacion}
Reglas largas suelen tener alta confianza, pero son difíciles de explicar.
\end{block}

\begin{columns}[T]
\column{0.48\textwidth}
\begin{lstlisting}[style=rcode]
table(size(rules_type))

rules_short <- rules_type[
  size(rules_type) <= 3
]
length(rules_short)
\end{lstlisting}

\column{0.48\textwidth}
\begin{lstlisting}[style=out]
  2   3   4   5   6   7   8   9  10
  3  29  95 155 150  92  37   9   1

# size <= 3
[1] 32
\end{lstlisting}
\end{columns}

\begin{block}{Lectura}
\texttt{size(rule)} cuenta LHS + RHS. Si RHS tiene un item, \texttt{size <= 3} equivale aproximadamente a LHS de maximo 2 items.
\end{block}
\end{frame}

%------------------------------------------------
\section{Redundancia}
\begin{frame}[fragile]{Filtro 6: reglas redundantes}
\begin{block}{Definicion operacional}
Una regla es redundante si existe una regla mas general con igual o mayor confianza.
\end{block}

\begin{columns}[T]
\column{0.52\textwidth}
\small
Ejemplo:
\[
\{eggs, fins\} \Rightarrow \{type=fish\}
\]
\[
\{eggs, aquatic, fins\} \Rightarrow \{type=fish\}
\]
Ambas tienen soporte 0.1287 y confianza 1. La segunda agrega \texttt{aquatic}, pero no mejora la regla.

\column{0.44\textwidth}
\begin{lstlisting}[style=rcode]
rules_nr <- rules_type[
  !is.redundant(rules_type)
]

rules_type
rules_nr
\end{lstlisting}
\begin{lstlisting}[style=out]
set of 571 rules
set of 29 rules
\end{lstlisting}
\end{columns}
\end{frame}

%------------------------------------------------
\begin{frame}[fragile]{Efecto del filtro de redundancia}
\begin{columns}[T]
\column{0.45\textwidth}
\begin{lstlisting}[style=rcode]
table(size(rules_type))
table(size(rules_nr))
\end{lstlisting}

\column{0.51\textwidth}
\begin{lstlisting}[style=out]
# Antes
  2   3   4   5   6   7   8   9  10
  3  29  95 155 150  92  37   9   1

# Despues
 2  3  4  5
 3 13 10  3
\end{lstlisting}
\end{columns}

\begin{center}
\begin{tikzpicture}
\begin{axis}[
  ybar, bar width=18pt,
  width=0.72\textwidth, height=0.32\textheight,
  ylabel={numero de reglas},
  symbolic x coords={Antes,Despues},
  xtick=data,
  ymin=0, ymax=620,
  nodes near coords,
  nodes near coords align={vertical},
]
\addplot coordinates {(Antes,571) (Despues,29)};
\end{axis}
\end{tikzpicture}
\end{center}
\end{frame}

%------------------------------------------------
\begin{frame}[fragile]{Reglas no redundantes: ejemplos interpretables}
\begin{lstlisting}[style=rcode]
inspect(rules_nr)
\end{lstlisting}

\vspace{0.1cm}
\scriptsize
\begin{tabular}{llrrrr}
\toprule
LHS & RHS & support & confidence & lift & count \\
\midrule
\{feathers\} & \{type=bird\} & 0.1980 & 1.0000 & 5.050 & 20 \\
\{milk\} & \{type=mammal\} & 0.4059 & 1.0000 & 2.463 & 41 \\
\{hair\} & \{type=mammal\} & 0.3861 & 0.9070 & 2.234 & 39 \\
\{eggs, fins\} & \{type=fish\} & 0.1287 & 1.0000 & 7.769 & 13 \\
\{fins, tail\} & \{type=fish\} & 0.1287 & 0.8125 & 6.312 & 13 \\
\{airborne, tail\} & \{type=bird\} & 0.1584 & 0.8889 & 4.489 & 16 \\
\{airborne, backbone\} & \{type=bird\} & 0.1584 & 0.8889 & 4.489 & 16 \\
\{hair, catsize\} & \{type=mammal\} & 0.2970 & 1.0000 & 2.463 & 30 \\
\bottomrule
\end{tabular}

\begin{alertblock}{Recordar}
Despues del filtro, las reglas son pocas y semanticas: plumas predicen aves, leche predice mamiferos, huevos+aletas predice peces.
\end{alertblock}
\end{frame}

%------------------------------------------------
\section{Medidas adicionales}
\begin{frame}[fragile]{Filtro 7: medidas adicionales con \texttt{interestMeasure}}
\begin{block}{Por que usar mas metricas?}
Soporte, confianza y lift no siempre capturan asociacion simetrica o poder discriminativo. \texttt{arules} permite calcular otras medidas, como \texttt{phi} y \texttt{gini}.
\end{block}

\begin{columns}[T]
\column{0.50\textwidth}
\begin{lstlisting}[style=rcode]
quality(rules) <- cbind(
  quality(rules),
  interestMeasure(
    rules,
    measure = c("phi", "gini"),
    trans = trans))

rules |>
  head(by = "phi") |>
  inspect()
\end{lstlisting}

\column{0.46\textwidth}
\scriptsize
\begin{tabular}{llrrrrr}
\toprule
LHS & RHS & supp. & conf. & lift & phi & gini \\
\midrule
\{eggs,fins\} & \{type=fish\} & 0.1287 & 1 & 7.769 & 1 & 0.224 \\
\{eggs,aquatic,fins\} & \{type=fish\} & 0.1287 & 1 & 7.769 & 1 & 0.224 \\
\{eggs,toothed,fins\} & \{type=fish\} & 0.1287 & 1 & 7.769 & 1 & 0.224 \\
\bottomrule
\end{tabular}
\end{columns}
\end{frame}

%------------------------------------------------
\begin{frame}[fragile]{Ejemplo de filtro combinado}
\begin{lstlisting}[style=rcode]
rules_filtered <- subset(
  rules,
  subset = support >= 0.05 &
           confidence >= 0.9 &
           lift > 2 &
           phi > 0.5
)

rules_filtered <- sort(rules_filtered, by = "lift")
inspect(head(rules_filtered, n = 10))
\end{lstlisting}

\begin{block}{Interpretacion}
Este filtro busca reglas con evidencia minima, alta predictividad, asociacion no trivial y correlacion fuerte entre LHS y RHS.
\end{block}

\begin{alertblock}{Recomendacion}
No existe un filtro universal. El filtro debe responder a la pregunta de analisis: exploracion, recomendacion, clasificacion o explicacion.
\end{alertblock}
\end{frame}

%------------------------------------------------
\section{Items negativos y codificacion}
\begin{frame}[fragile]{Codificacion: incluir items negativos puede explotar las reglas}
\begin{columns}[T]
\column{0.51\textwidth}
\begin{lstlisting}[style=rcode]
Zoo_factors <- Zoo_has_legs |>
  mutate(across(where(is.logical), factor))

trans_factors <- transactions(Zoo_factors)
trans_factors

r_neg <- apriori(trans_factors)
r_neg
print(object.size(r_neg), unit = "Mb")
\end{lstlisting}

\column{0.45\textwidth}
\begin{lstlisting}[style=out]
transactions in sparse format with
 101 transactions
 39 items

writing ... [1517191 rule(s)]

set of 1517191 rules
110.2 Mb
\end{lstlisting}
\end{columns}

\begin{alertblock}{Interpretacion}
Convertir TRUE/FALSE a factores crea items como \texttt{milk=FALSE}, \texttt{fins=FALSE}. Esto puede revelar ausencias informativas, pero tambien genera muchas reglas obvias.
\end{alertblock}
\end{frame}

%------------------------------------------------
\begin{frame}[fragile]{Ejemplo: reglas con items negativos}
\begin{lstlisting}[style=rcode]
r_neg |>
  head(n = 10, by = "lift") |>
  inspect()
\end{lstlisting}

\scriptsize
\begin{tabular}{llrrrr}
\toprule
LHS & RHS & support & confidence & lift & count \\
\midrule
\{breathes=FALSE, fins=TRUE\} & \{type=fish\} & 0.1287 & 1 & 7.769 & 13 \\
\{eggs=TRUE, fins=TRUE\} & \{type=fish\} & 0.1287 & 1 & 7.769 & 13 \\
\{milk=FALSE, fins=TRUE\} & \{type=fish\} & 0.1287 & 1 & 7.769 & 13 \\
\{breathes=FALSE, fins=TRUE, legs=FALSE\} & \{type=fish\} & 0.1287 & 1 & 7.769 & 13 \\
\{aquatic=TRUE, breathes=FALSE, fins=TRUE\} & \{type=fish\} & 0.1287 & 1 & 7.769 & 13 \\
\bottomrule
\end{tabular}

\begin{block}{Cuando usar items negativos}
Usarlos selectivamente cuando la ausencia tenga sentido: por ejemplo, \texttt{milk=FALSE} puede ayudar a diferenciar peces de mamiferos, pero incluir todas las ausencias puede saturar el analisis.
\end{block}
\end{frame}

%------------------------------------------------
\section{Representaciones compactas}
\begin{frame}[fragile]{Itemsets cerrados y maximales}
\begin{block}{Motivacion}
Antes de mirar reglas, se puede reducir el espacio mirando representaciones compactas de itemsets frecuentes.
\end{block}

\begin{columns}[T]
\column{0.52\textwidth}
\begin{lstlisting}[style=rcode]
its <- apriori(trans,
  parameter = list(target = "frequent",
                   support = 0.05))

its_closed <- its[is.closed(its)]
its_max <- its[is.maximal(its)]

c(frequent = length(its),
  closed = length(its_closed),
  maximal = length(its_max))
\end{lstlisting}

\column{0.44\textwidth}
\begin{lstlisting}[style=out]
frequent  closed maximal
    2537     230      22
\end{lstlisting}

\begin{block}{Lectura}
\texttt{closed} conserva soporte exacto. \texttt{maximal} conserva solo patrones mas especificos, pero pierde informacion de soporte de subconjuntos.
\end{block}
\end{columns}
\end{frame}

%------------------------------------------------
\begin{frame}[fragile]{Ejemplos de itemsets maximales y cerrados}
\begin{columns}[T]
\column{0.47\textwidth}
\begin{lstlisting}[style=rcode]
its_max |>
  head(by = "support") |>
  inspect()
\end{lstlisting}
\scriptsize
\begin{tabular}{lr}
\toprule
Itemset maximal & count \\
\midrule
\{hair, milk, predator, toothed, backbone, breathes, legs, tail, catsize, type=mammal\} & 13 \\
\{eggs, aquatic, predator, toothed, backbone, fins, tail, type=fish\} & 9 \\
\bottomrule
\end{tabular}

%\column{0.49\textwidth}
%\begin{lstlisting}[style=rcode]
%its_closed |>
%  head(by = "support") |>
%  inspect()
%\end{lstlisting}
%\scriptsize
%\begin{tabular}{lrr}
%\toprule
%Itemset cerrado & support & count \\
%\midrule
%\{backbone\} & 0.8218 & 83 \\
%\{breathes\} & 0.7921 & 80 \\
%\{legs\} & 0.7723 & 78 \\
%\{tail\} & 0.7426 & 75 \\
%\{backbone, tail\} & 0.7327 & 74 \\
%\{breathes, legs\} & 0.7228 & 73 \\
%\bottomrule
%\end{tabular}
\end{columns}
\end{frame}

%\begin{frame}[fragile]{Flujo recomendado para los alumnos}
%\begin{enumerate}
%  \item Preparar transacciones: revisar codificacion, items raros y frecuentes.
%  \item Minar reglas con soporte conservador; bajar soporte solo si se justifica.
%  \item Ordenar por lift, pero reportar tambien support, confidence, coverage y count.
%  \item Usar \texttt{appearance} si se busca una clase o consecuente especifico.
%  \item Filtrar reglas largas: empezar con LHS de 1-2 items.
%  \item Eliminar redundantes con \texttt{is.redundant()}.
%  \item Visualizar: item frequency $\rightarrow$ scatter $\rightarrow$ matrix/graph para top reglas.
%  \item Interpretar semanticamente: descartar reglas obvias o artefactos de codificacion.
%\end{enumerate}
%\end{frame}
%
%%------------------------------------------------
%\begin{frame}[fragile]{Mini-guia de codigo final}
%\begin{lstlisting}[style=rcode]
%data(Zoo, package = "mlbench")
%Zoo_has_legs <- Zoo |> mutate(legs = legs > 0)
%trans <- transactions(Zoo_has_legs)
%
%# Reglas orientadas a clasificacion de tipo animal
%type <- grep("type=", itemLabels(trans), value = TRUE)
%rules_type <- apriori(trans, appearance = list(rhs = type))
%
%# Medidas adicionales
%quality(rules_type) <- cbind(
%  quality(rules_type),
%  interestMeasure(rules_type, c("phi", "gini"), trans = trans))
%
%# Filtro interpretable
%rules_final <- rules_type[!is.redundant(rules_type)]
%rules_final <- subset(rules_final,
%                      subset = lift > 2 & size(rules_final) <= 5)
%rules_final <- sort(rules_final, by = "lift")
%inspect(rules_final)
%
%# Visualizacion
%plot(rules_final)
%plot(rules_final, method = "graph")
%\end{lstlisting}
%\end{frame}



%%------------------------------------------------
%\section{Mini taller}
%
%\begin{frame}{Taller en clase}
%\begin{enumerate}
%  \item Minar reglas en \texttt{Groceries} con tres valores de soporte.
%  \item Filtrar reglas con \texttt{lift > 1.5}, \texttt{count > 50} y \texttt{confidence > 0.5}.
%  \item Eliminar reglas redundantes.
%  \item Crear un scatterplot y una matriz agrupada.
%  \item Escoger 3 reglas y justificar si son accionables o no.
%\end{enumerate}
%\end{frame}
%
%%------------------------------------------------
%\begin{frame}[fragile]{Código base del taller}
%\begin{lstlisting}
%data("Groceries")
%
%rules <- apriori(
%  Groceries,
%  parameter = list(support = 0.005, confidence = 0.3, minlen = 2)
%)
%
%quality(rules) <- cbind(
%  quality(rules),
%  interestMeasure(rules,
%                  measure = c("leverage", "conviction", "phi"),
%                  transactions = Groceries)
%)
%
%rules_clean <- subset(rules, lift > 1.5 & confidence > 0.5 & count > 50)
%rules_clean <- rules_clean[!is.redundant(rules_clean)]
%rules_clean <- sort(rules_clean, by = "lift")
%
%inspect(head(rules_clean, 10))
%plot(rules_clean)
%plot(head(rules_clean, 50), method = "graph")
%\end{lstlisting}
%\end{frame}

%------------------------------------------------
\begin{frame}{Checklist para reportar reglas}
\begin{enumerate}
  \item ¿Cómo se codificaron o discretizaron los datos?
  \item ¿Cuáles fueron los umbrales de soporte, confianza y largo de regla?
  \item ¿Qué filtros posteriores se aplicaron?
  \item ¿Se eliminaron reglas redundantes?
  \item ¿Se revisaron conteos absolutos?
  \item ¿Se visualizaron patrones globales y no solo reglas individuales?
  \item ¿La interpretación evita causalidad no justificada?
\end{enumerate}
\end{frame}

%------------------------------------------------
\begin{frame}{Resumen}
\begin{itemize}
  \item El valor del análisis de asociación no está solo en minar reglas, sino en filtrarlas e interpretarlas.
  \item Lift, leverage, conviction, count y significancia ayudan a priorizar reglas.
  \item Las visualizaciones permiten detectar estructura, redundancia y grupos de reglas.
  \item La discretización y las jerarquías cambian el espacio de patrones disponibles.
  \item Un buen informe explica el criterio de selección.
\end{itemize}
\end{frame}

%------------------------------------------------
\begin{frame}{Referencias }
\footnotesize
\begin{itemize}
  \item Michael Hahsler. \textit{An R Companion for Introduction to Data Mining}. Capítulo 5: Association Analysis: Basic Concepts.\\
  \url{https://mhahsler.github.io/Introduction_to_Data_Mining_R_Examples/book/association-analysis-basic-concepts.html}
  \item Michael Hahsler. \textit{An R Companion for Introduction to Data Mining}. Capítulo 6: Association Analysis: Advanced Concepts.\\
  \url{https://mhahsler.github.io/Introduction_to_Data_Mining_R_Examples/book/association-analysis-advanced-concepts.html}
  \item Hahsler, M., Grün, B., Hornik, K. (2005). \textit{arules: Mining Association Rules and Frequent Itemsets}.
  \item Hahsler, M. (2017). \textit{arulesViz: Interactive Visualization of Association Rules with R}.
\end{itemize}
\end{frame}

\end{document}
